% Standalone problem document, generated from the adjacent README.md.
% Compile with XeLaTeX. Mathematical content is maintained in Markdown.
\documentclass[11pt,a4paper]{article}
\usepackage[fontsize=10.5pt]{scrextend}
\usepackage[margin=22mm,headheight=15pt,headsep=9mm,footskip=11mm]{geometry}
\usepackage{fontspec}
\setmainfont{texgyrepagella}[Extension=.otf,UprightFont=*-regular,BoldFont=*-bold,ItalicFont=*-italic,BoldItalicFont=*-bolditalic]
\setsansfont{texgyreheros}[Extension=.otf,UprightFont=*-regular,BoldFont=*-bold,ItalicFont=*-italic,BoldItalicFont=*-bolditalic]
\setmonofont{texgyrecursor}[Extension=.otf,UprightFont=*-regular,BoldFont=*-bold,ItalicFont=*-italic,BoldItalicFont=*-bolditalic,Scale=MatchLowercase]
\usepackage{amsmath,amssymb}
\usepackage{unicode-math}
\setmathfont{texgyrepagella-math.otf}
\usepackage{xcolor,microtype,enumitem,needspace,fancyhdr}
\usepackage{bookmark}
\definecolor{ink}{HTML}{17344B}
\definecolor{accent}{HTML}{197D80}
\definecolor{muted}{HTML}{596773}
\hypersetup{colorlinks=true,linkcolor=accent,urlcolor=accent,pdftitle={TR-29 - Exact
partially symmetric rank of products of generalized W
tensors},pdfauthor={Open Problems in Numerical Linear Algebra}}
\urlstyle{same}
\setlength{\parindent}{0pt}
\setlength{\parskip}{5pt plus 1pt minus 1pt}
\linespread{1.04}
\setlength{\emergencystretch}{3em}
\widowpenalty=10000
\clubpenalty=10000
\displaywidowpenalty=10000
\allowdisplaybreaks[1]
\setlist{nosep,leftmargin=1.5em}
\providecommand{\tightlist}{\setlength{\itemsep}{0pt}\setlength{\parskip}{0pt}}
\providecommand{\pandocbounded}[1]{#1}
\pagestyle{fancy}
\fancyhf{}
\fancyhead[L]{\sffamily\footnotesize\color{muted}OPEN PROBLEMS IN NUMERICAL LINEAR ALGEBRA}
\fancyhead[R]{\sffamily\bfseries\footnotesize\color{accent}TR-29}
\fancyfoot[L]{\parbox[b]{0.9\textwidth}{\sffamily\fontsize{8}{10}\selectfont\color{muted}Editorial ratings; a bounded search cannot certify that a problem remains unsolved.\\Literature check: 2026-09-10}}
\fancyfoot[R]{\sffamily\footnotesize\color{muted}\thepage}
\renewcommand{\headrulewidth}{0.3pt}
\renewcommand{\headrule}{\hbox to\headwidth{\color{accent}\leaders\hrule height \headrulewidth\hfill}}
\makeatletter
\renewcommand{\section}{\Needspace{5\baselineskip}\@startsection{section}{1}{0pt}{12pt plus 2pt minus 2pt}{4pt}{\sffamily\bfseries\large\color{ink}}}
\renewcommand{\subsection}{\Needspace{4\baselineskip}\@startsection{subsection}{2}{0pt}{9pt plus 1pt minus 1pt}{3pt}{\sffamily\bfseries\normalsize\color{ink}}}
\makeatother
\setcounter{secnumdepth}{0}
\begin{document}
{\sffamily\small\color{muted}Tensor computations\par}
\vspace{3pt}
{\raggedright\hyphenpenalty=10000\exhyphenpenalty=10000\sffamily\bfseries\fontsize{21}{24}\selectfont\color{ink}Exact
partially symmetric rank of products of generalized W tensors\par}
\vspace{7pt}
{\sffamily\small\textbf{Difficulty:} challenging\quad\textbf{Importance:} interesting
to specialist\par}
{\sffamily\small\bfseries\color{accent}Status: Partially resolved\par}
\vspace{4pt}
{\color{accent}\hrule height 0.8pt}
\vspace{6pt}
\textbf{Rating rationale:} Matching constructive upper bounds by sharp
lower bounds for arbitrary products requires substantial new
decomposition arguments. The exact partially symmetric rank of this
particular family chiefly interests specialists in tensor rank and
multipartite entanglement.

\subsection{Problem statement}\label{problem-statement}

For an integer \(d\geq3\), put \[
W_d=\sum_{j=1}^{d}
e_1^{\otimes(j-1)}\otimes e_2\otimes
e_1^{\otimes(d-j)}\in\operatorname{Sym}^d(\mathbb C^2).
\] Given any \(k\geq2\) and \(d_1,\ldots,d_k\geq3\), set
\(T_{\mathbf d}=W_{d_1}\otimes\cdots\otimes W_{d_k}\). Determine the
exact smallest integer \(s\) such that \[
T_{\mathbf d}=\sum_{\ell=1}^{s}
(v_{1,\ell})^{\otimes d_1}\otimes\cdots\otimes
(v_{k,\ell})^{\otimes d_k},
\qquad v_{j,\ell}\in\mathbb C^2.
\] This minimum is the partially symmetric rank with the \(k\) blocks of
orders \(d_1,\ldots,d_k\) fixed. It is not the border rank and not the
rank obtained after grouping modes from different blocks. The harmless
nonzero scalar relating the displayed \(W_d\) to the polynomial
\(x^{d-1}y\) can be absorbed into a decomposition.

\subsection{Why it matters}\label{why-it-matters}

These explicit tensors are benchmarks for the difference between exact,
border, and symmetry-constrained tensor decompositions. They already
demonstrate that decomposing two copies together can cost fewer terms
than multiplying their individual ranks. Determining their exact ranks
would quantify that saving beyond small examples.

\subsection{References and status
check}\label{references-and-status-check}

\begin{enumerate}
\def\labelenumi{\arabic{enumi}.}
\tightlist
\item
  A. Oneto and E. Ventura, \emph{Ranks of tensors: geometry and
  applications}, \href{https://doi.org/10.1007/s40574-025-00472-9}{2025
  survey}, Question 12, following Theorems 4.4--4.5.
\item
  E. Ballico, A. Bernardi, M. Christandl, and F. Gesmundo, \emph{On the
  partially symmetric rank of tensor products of W-states and other
  symmetric tensors},
  \href{https://arxiv.org/abs/1803.01623}{arXiv:1803.01623}, §§3--4;
  Rendiconti Lincei \textbf{30} (2019), 93--124.
\item
  S. Canino, A. Casarotti, and P. Santarsiero, \emph{A new bound on the
  rank of tensor product of W-states},
  \href{https://arxiv.org/html/2512.05828v1}{arXiv:2512.05828v1},
  Theorem 1.1 and its following discussion.
\end{enumerate}

\subsubsection{Status check ---
2026-09-10}\label{status-check-2026-09-10}

Checked the 2025 source question and Canino--Casarotti--Santarsiero v1,
Theorem 1.1 and the ensuing sharpness discussion, with targeted
W-product-rank searches through 2026. The December 2025 paper gives the
upper bound \(2^{k-1}(d_1+\cdots+d_k-2k+2)\) and explicit
decompositions. Equality is known for \(k=2,d_1=d_2=3\), where the rank
is eight, so a nontrivial parameter case is resolved. The construction
alone does not prove minimality for other tuples. No general exact
formula was located.
\end{document}
